\section{Syntactic categories}
\label{sec:b-syntax}

The formalization distinguishes three syntactic categories: basic types and two
representations of B terms.
The concrete, first-order representation closely follows the syntax obtained
from Atelier B's proof-obligation files and thereby provides a deep embedding of the B
language.
The abstract representation uses parametric higher-order abstract syntax
(PHOAS), as motivated in \Cref{sec:phoas}, and serves as the internal syntax
for semantic reasoning.

\subsection{Concrete terms}
\label{subsec:b-concrete-terms}

\paragraph{Basic terms.}
We first define a syntactic category for concrete B terms, thereby giving a
first-order representation of a curated fragment of the B object language.
This deep embedding constitutes the concrete layer of the formalization in
Lean.

\begin{definition}[Concrete terms][B.Term][B/Syntax/Basic.html\#B.Term]
  \label{def:b-basic-syntax}
  The constructors of \leaninl{B.Term} are listed below, where \(\var{v}\) is a name in
  \leaninl{String}, \(n\) is an integer, \(b\) is a Boolean value, and
  \(\vec{\var{v}}\) is a list of names. Object variables are typeset in bold gray
  throughout.

  \begin{bnf}[colspec = {l@{}l@{\kern6pt}c@{\kern6pt}ll}]
    \leaninl{B.Term} $\,∋\, t, s$ ::=
    | $\var{v}$ && $\mathsf{int}\ n$ && $\mathsf{bool}\ b$ && $\ZintB$ && $\ZboolB$ // Atoms
    | $t \regnotation{\mapletB} s$ // Maplet
    | $t \regnotation{\addB} s$ && $t \regnotation{\subB} s$ && $t \regnotation{\mulB} s$ && $t \regnotation{\leB} s$ // Arithmetic
    | $t \regnotation{\andB} s$ && $\regnotation{\notB} t$ && $t \regnotation{\eqB} s$ && $t \inB s$ && $\forallB \vec{\var{v}} \inB t \cdot s$ // Logic
    | $\powB(t)$ && $t \prodB s$ && $t \regnotation{\cupB} s$ && $t \regnotation{\capB} s$ && $\regnotation{\bcomp{\vec{\var{v}} \inB t}{s}}$ // Set theory
    | $\cardB{t}$
    | $t \pfunB s$ && $\appB t\ s$ && $\regnotation{\lambdaB} \vec{\var{v}} \inB t \cdot\,s$ // Functions
  \end{bnf}
\end{definition}

\paragraph{Additional concrete terms.}

In terms of expressiveness, the grammar defined above is near-minimal: its
primitive signature has been reduced to a small set of constructors that
suffice to express a significant fragment of the B language. It is not strictly
minimal, since some constructs could be introduced as definitional extensions,
but treating them as primitive simplifies the subsequent encoding.

\begin{example}
  For instance, universal quantification
  \( \forallB \vec{\var{v}} \inB D \cdot P \) can be characterized by the following
  equality:
  \[
    (\lambdaB \vec{\var{v}} \inB D \cdot P) \eqB (\lambdaB \vec{\var{v}} \inB D \cdot \mathsf{bool}\ \mathsf{true}).
  \]
\end{example}

As the preceding example suggests, such reductions are justified semantically
rather than purely syntactically.
We nevertheless define additional B connectives as derived forms, as listed below.

\begin{definition}[Additional concrete terms][B.Term, String.toBinaryOp][B/Syntax/Extra.html, POGReader/POGReader.html\#String.toBinaryOp]
  \label{def:b-extra-syntax}
  The following B constructs are defined as derived forms in terms of the
  primitive syntax of \Cref{def:b-basic-syntax}. Let \(x, S, T, P, Q\) be B terms, and let
  \(\vec{\var{v}}\) be a list of concrete variable names.

  \textbf{Logical connectives:}
  \begin{align*}
    x \regnotation{\notinB} S                     & \triangleq \notB(x \inB S) \tag{negated membership}                                      \\
    P \regnotation{\orB} Q                        & \triangleq \notB((\notB P) \andB (\notB Q)) \tag{disjunction}                            \\
    P \regnotation{\impB} Q                       & \triangleq (\notB P) \orB Q \tag{implication}                                            \\
    P \regnotation{\iffB} Q                       & \triangleq (P \impB Q) \andB (Q \impB P) \tag{bi-implication}                            \\
    P \regnotation{\neqB} Q                       & \triangleq \notB(P \eqB Q) \tag{negated equality}                                        \\
    \regnotation{\existsB} \vec{\var{v}} \inB S \cdot P & \triangleq \notB(\forallB \vec{\var{v}} \inB S \cdot \notB P) \tag{existential quantification}
    \intertext{\textbf{Relations:}}
    S \bop{\leftrightarrow} T       & \triangleq \powB (S \prodB T) \tag{set of relations}
  \end{align*}
\end{definition}

\subsection{Basic types}
\label{subsec:b-type-system}
We also introduce a syntactic category for basic types.
Because it is grounded in set theory, the B mathematical language does not,
strictly speaking, include an intrinsic type system.
The B-Book~\cite{abrial1996bbook} introduces a relation associated with a
typing judgment, but the objects constrained by this relation belong to the
semantic account of the language rather than to a separate type-theoretic
stratum.
This relation is intended to ensure the well-definedness of terms and should be
understood as a membership constraint that underlies the generation of proof
obligations.
In this sense, the B language is formally untyped.

Nevertheless, a type system provides useful support for formalizing the
language, manipulating well-defined expressions, and defining objects that
cannot easily be represented in pure set theory.
We therefore introduce an auxiliary syntactic layer of objects, called
``B types'', with four base forms: integers, Booleans, (power) sets, and
(Cartesian) products.
Typing rules relating terms to types are introduced later; the present section
is purely syntactic.

\begin{definition}[B types][B.BType][B/Typing/Basic.html\#B.BType]
  \label{def:b-types}
              \begin{bnf}
    \leaninl{B.BType} \(\,∋\, \tau\) ::= \(\regnotation{\intB}\) && \(\regnotation{\boolB}\) && \(\regnotation{\setB} \tau\) && \(\tau \prodB \tau\)
  \end{bnf}
\end{definition}

\begin{remark}
  This set of types could easily be extended to model user-defined sets, which
  can be used as type parameters for functions and relations. For the moment,
  such sets, along with enumerations, are represented as sets of integers;
  enumerations, in particular, are finite and non-empty by definition, and this
  representation is adequate for the fragment we model.
  Moreover, this representation does not include propositions as a distinct
  syntactic category; we regard them as being on the same level as Booleans,
  although they are technically distinct syntactic categories in B.
\end{remark}

B types admit a natural carrier interpretation as concrete B terms.
Although the present section is purely syntactic, this embedding is
straightforward, and we give its definition for completeness.
Since the embedding lifts B types to the term level, we denote the concrete term
associated with a B type \(\tau\) by \(\lift{\tau}\).

\begin{definition}[B types as terms][B.BType.toTerm][B/Typing/Basic.html\#B.BType.toTerm]
  \label{def:btype-to-term}
    The embedding of B types into concrete B terms is defined by structural
  recursion as follows:
  \begin{equation*}
    \regnotation{\lift{\intB}} \triangleq \ZintB,\quad
    \lift{\boolB} \triangleq \ZboolB,\quad
    \lift{(\setB\tau)} \triangleq \powB(\lift{\tau}),\quad
    \lift{(α \prodB β)} \triangleq \powB(\lift{α} \prodB \lift{β})
  \end{equation*}
\end{definition}

\subsection{Additional terms}

At this stage, the primitive syntax covers only a small subset of the B
language; nevertheless, richer compound constructs can be introduced as
definitional extensions.
The derived forms in \Cref{def:b-extra-syntax} do not bind new variables.
More complex constructs, however, can be expressed by bounded comprehensions, as
illustrated by the following example.

\begin{example}
  Although intersection \( t \capB s \) is available as a primitive
  (\Cref{def:b-basic-syntax}), it can equally be expressed as the bounded comprehension
  \[
    \bcomp{\var{v} \inB t}{\var{v} \inB s},
  \]
  which illustrates the pattern: the definition captures the intended semantics of
  intersection and, in particular, yields an associative and commutative operation.
\end{example}

Defining such constructs requires the ability to bind variables, which is not
directly available at the concrete syntax level.
At the implementation level, however, binding is available when parsing and
decoding proof obligations.
We can therefore define many additional constructs by relying on the parser,
its binding infrastructure, and the intended semantic interpretation of those
constructs.

The following definitions are based on the specification of the B language in the
B-Book~\cite{abrial1996bbook}.
Some symbols are annotated with a B type; in such cases, the parser ensures
that the symbol is used only with terms of the specified type.
Although the typing rules are introduced only later, in \Cref{sec:b-typing},
the parser is designed to fail when a type does not match the expected shape.
These B types are therefore used as parsing-time well-formedness constraints on
the syntax of terms.

\paragraph{Set-theoretic constructs.}
We add support for further set-theoretic constructs.
One of the most basic such objects is the empty set, defined as a set with no
elements.
Since set comprehension in the B language is bounded, ``the'' empty set is not
unique and must be defined for each element type.
This type dependence is not visible in the surface syntax: in a B machine, one
can write \(\varnothing\) or \(\{\}\), since well-defined B expressions
implicitly constrain the type of the empty set.
We also provide definitions for the subset relation, set difference, and the
non-empty powerset.

\begin{definition}[Set-theoretic constructs and operations][B.Builtins, toBinaryOp, toUnaryOp][
    /POGReader/Builtins.html,
    POGReader/POGReader.html\#String.toBinaryOp,
    POGReader/POGReader.html\#String.toUnaryOp]
  \label{def:extra-parsed-syntax}
  Let \(S\) and \(T\) be B terms, and let \(\tau\) be a B type fixed by the parser, as
  annotated on the right of each equation.
  \begin{align*}
    \varnothing_{\tau}  & \;\triangleq\; \bcomp{\var{v} \inB \lift{\alpha}}{\mathsf{bool}\ \mathsf{false}} \tag{\(\tau = \setB(\alpha)\)}            \\
    S \regnotation{\subseteqB} T      & \;\triangleq\; \forallB \var{v} \inB S \cdot \var{v} \inB T                                           \\
    S \regnotation{\notsubseteqB} T   & \;\triangleq\; \notB(S \subseteqB T)                                                      \\
    S \regnotation{\ssubB} T          & \;\triangleq\; \bcomp{\var{v} \inB S}{\var{v} \notinB T}                                          \\
    {\powB_1}_{\tau}(S) & \;\triangleq\; \bcomp{\var{v} \inB \powB(S)}{\existsB \var{x} \inB \lift{\alpha} \cdot \var{x} \inB \var{v}} \tag{\(\tau = \setB(\setB(\alpha))\)}
  \end{align*}
\end{definition}

We also add finite sets, since they are ubiquitous in B specifications and are
not provided as primitive constructs in the set-theoretic core.
A subset \(S\) of a set \(A\) is finite when there exists an injection from
\(S\) into a bounded interval of integers, typically one whose lower bound is
\(0\).

\begin{definition}[Finite subsets][B.Term.finite][POGReader/Builtins.html\#B.Term.finite]
  \label{def:b-syntax-finite}
  Let \(S\) be a B term, and let \(\tau\) be a B type;
  the total-injection space \(\injtfunB\) used below is introduced in \Cref{def:b-syntax-functions}.
  \begin{align*}
    \mathsf{FIN}(S)            & \;\triangleq\; \bcomp{\var{X} \inB \powB(S)}{\existsB \var{n} \inB \ZintB \cdot \existsB \var{f} \inB S \injtfunB \ZintB \cdot \forallB \var{x}, \var{i} \inB S \prodB \ZintB \cdot \\
                               & \hspace{1cm}(\var{x} \mapletB \var{i} \inB \var{f}) \impB ((\mathsf{int}\ 0 \leB \var{i}) \andB (\var{i} \leB \var{n}))}                                                            \\
    {\mathsf{FIN}_1}_{\tau}(S) & \;\triangleq\; \bcomp{\var{X} \inB \mathsf{FIN}(S)}{\var{X} \neqB \varnothing_{\tau}} \tag{\(\tau = \setB(\alpha)\)}
  \end{align*}
\end{definition}

\paragraph{Arithmetic constructs.}
We add support for basic arithmetic constructs, including the set of natural
numbers as a subset of integers\footnote{This marks another deviation from
  pure set theory: the natural numbers are introduced here as a distinguished
  subset of the integers.}, integer intervals, and comparison operators.

\begin{definition}[Arithmetic constructs][B.Builtins, toBinaryOp][
    POGReader/Builtins.html,
    POGReader/POGReader.html\#String.toBinaryOp]
  \label{def:b-syntax-arithmetic}
          Let \(i, j\) be B terms.
  \begin{align*}
    \mathbb{N}       & \;\triangleq\; \bcomp{\var{n} \inB \ZintB}{\mathsf{int}\ 0 \leB \mathsf{int}\ \var{n}} \\
    \regnotation{\intrangeB{i}{j}} & \;\triangleq\; \bcomp{\var{n} \inB \ZintB}{i \leB \var{n} \andB \var{n} \leB j}              \\
    i \regnotation{\geB} j         & \;\triangleq\; j \leB i                                                          \\
    i \regnotation{\gtB} j         & \;\triangleq\; \notB(i \leB j)                                                   \\
    i \regnotation{\ltB} j         & \;\triangleq\; \notB(j \leB i)
  \end{align*}
\end{definition}

\paragraph{Relations.}
Relations form the largest class of derived constructs, since they are heavily
used in B specifications.
We therefore introduce the identity relation, domain and range operators with
their associated restriction and subtraction operations, projections,
overriding, direct product, inverse relation, and direct image.

The identity relation is the graph of the identity map on a set.
\begin{definition}[Identity relation/function][toUnaryOp][
    POGReader/POGReader.html\#String.toUnaryOp]
  \label{def:b-syntax-id}
  Let \(S\) be a B term.
  \[
    \mathsf{id}(S) \;\triangleq\; \lambdaB \var{v} \inB S \cdot \var{v}
  \]
\end{definition}

Given a relation \(R\) and a set \(F\), the domain restriction of \(R\) to
\(F\) is the subrelation obtained by retaining only the pairs whose first
component belongs to \(F\). The domain subtraction of \(R\) by \(F\) is the
subrelation obtained by retaining only the pairs whose first component does not
belong to \(F\).
\begin{definition}[Domain restriction and subtraction][B.Term.domRestriction, B.Term.domSubtraction][
    POGReader/Builtins.html\#B.Term.domRestriction,
    POGReader/Builtins.html\#B.Term.domSubtraction]
  \label{def:b-syntax-domRestrSub}
      Let \(F\) and \(R\) be B terms.
  \[
    F \regnotation{\domRestrB} R \;\triangleq\; \bcomp{\var{x}, \var{y} \inB R}{\var{x} \inB F} \quad\text{and}\quad
    F \regnotation{\domSubtrB} R \;\triangleq\; \bcomp{\var{x}, \var{y} \inB R}{\var{x} \notinB F}
  \]
\end{definition}

Given a relation \(R\) and a set \(F\), the range restriction of \(R\) to \(F\)
is the subrelation obtained by retaining only the pairs whose second component
belongs to \(F\). The range subtraction of \(R\) by \(F\) is the subrelation
obtained by retaining only the pairs whose second component does not belong to
\(F\).
\begin{definition}[Range restriction and subtraction][B.Term.ranRestriction, B.Term.ranSubtraction][
    POGReader/Builtins.html\#B.Term.ranRestriction,
    POGReader/Builtins.html\#B.Term.ranSubtraction]
  \label{def:b-syntax-ranRestrSub}
      Let \(F\) and \(R\) be B terms.
  \[
    R \regnotation{\ranRestrB} F \;\triangleq\; \bcomp{\var{x}, \var{y} \inB R}{\var{y} \inB F} \quad\text{and}\quad
    R \regnotation{\ranSubtrB} F \;\triangleq\; \bcomp{\var{x}, \var{y} \inB R}{\var{y} \notinB F}
  \]
\end{definition}

The domain and range of a relation \(R\) are the sets of first and second
components, respectively, of the pairs in \(R\).
\begin{definition}[Domain and range of a relation][B.Term.dom, B.Term.ran][
    POGReader/Builtins.html\#B.Term.dom,
    POGReader/Builtins.html\#B.Term.ran]
  \label{def:b-syntax-dom-ran}
      Let \(R\) be a B term, and let \(\tau\) be a B type fixed by the parser, as annotated on the right.
  \begin{align*}
    \domB[\tau](R) & \;\triangleq\; \bcomp{\var{x} \inB \lift{\alpha}}{\existsB \var{y} \inB \lift{\beta} \cdot (\var{x} \mapletB \var{y}) \inB R} \tag{\(\tau = \setB(\alpha \prodB \beta)\)} \\
    \ranB[\tau](R) & \;\triangleq\; \bcomp{\var{y} \inB \lift{\beta}}{\existsB \var{x} \inB \lift{\alpha} \cdot (\var{x} \mapletB \var{y}) \inB R} \tag{\(\tau = \setB(\alpha \prodB \beta)\)}
  \end{align*}
\end{definition}

Projections are defined as coordinate relations that associate each pair in a
Cartesian product with its first and second components, respectively.
\begin{definition}[Projections][toBinaryOp][
    POGReader/POGReader.html\#String.toBinaryOp]
  Let \(Q\) and \(R\) be B terms.
  \begin{align*}
    \mathsf{prj}_1(Q, R) & \;\triangleq\; \bcomp{\var{x}, \var{y}, \var{z} \inB Q \prodB R \prodB Q}{\var{z} \eqB \var{x}} \\
    \mathsf{prj}_2(Q, R) & \;\triangleq\; \bcomp{\var{x}, \var{y}, \var{z} \inB Q \prodB R \prodB R}{\var{z} \eqB \var{y}}
  \end{align*}
\end{definition}

\begin{note}
  Recall that the Cartesian product \(\prodB\) is left associative in B, so that
  \(Q \prodB R \prodB Q\) is parsed as \((Q \prodB R) \prodB Q\).
\end{note}

The overriding operation on relations contains all the pairs in the overriding
relation \(R\), together with all pairs in the overridden relation \(Q\) whose
first component does not occur as the first component of any pair in \(R\).

\begin{definition}[Overloading][B.Term.overload][
    POGReader/Builtins.html\#B.Term.overload]
  Let \(Q\) and \(R\) be B terms, and let \(\tau\) be a B type fixed by the parser.
  \[
    Q \overloadB[\tau] R \;\triangleq\; \bcomp{\var{x}, \var{y} \inB \lift{\alpha}\! \prodB \lift{\beta}}{\left(\left(\var{x} \mapletB \var{y} \inB Q\right) \andB \notB(\var{x} \inB \domB[\tau](R))\right) \orB (\var{x} \mapletB \var{y} \inB R)}
  \]
  where \(\tau = \setB(\alpha \prodB \beta)\).
\end{definition}

The inverse relation of a relation \(R\) is obtained by swapping the first and
second components of each pair in \(R\).
\begin{definition}[Inverse relation][toUnaryOp][
    POGReader/POGReader.html\#String.toUnaryOp]
  \label{def:b-syntax-inv-rel}
  Let \(R\) be a B term, and let \(\tau\) be a B type fixed by the parser.
  \[
    R^{-1}_{\tau} \;\triangleq\; \bcomp{\var{x}, \var{y} \inB \lift{\beta} \prodB \lift{\alpha}}{\var{y} \mapletB \var{x} \inB R} \quad (\tau = \setB(\alpha \prodB \beta))
  \]
\end{definition}

The direct image of a set \(S\) under a relation \(R\) is the set of elements
that are related by \(R\) to some element of \(S\).
\begin{definition}[Direct image][String.toBinaryOp][
    POGReader/POGReader.html\#String.toBinaryOp]
  \label{def:b-syntax-dir-image}
  Let \(R\) and \(S\) be B terms, and let \(\tau\) be a B type fixed by the parser.
  \[
    R[S]_{\tau} \;\triangleq\; \bcomp{\var{y} \inB \lift{\alpha}}{\existsB \var{x} \inB S \cdot \var{x} \mapletB \var{y} \inB R} \tag{\(\tau = \setB(\alpha)\)}
  \]
\end{definition}

The direct product of two relations \(Q\) and \(R\) is the set of pairs \(x \mapletB (y \mapletB z)\) such that
\(x\) is related to \(y\) by \(Q\) and to \(z\) by \(R\).
\begin{definition}[Direct product][toBinaryOp][
    POGReader/POGReader.html\#String.toBinaryOp]
  \label{def:b-syntax-dirprod}
  Let \(Q\) and \(R\) be B terms, and let \(\tau\) be a B type fixed by the parser.
  \begin{fitalign}
    Q \dirprodB[\tau] R & \;\triangleq\; \bcomp{
    \var{x}, \var{w} \inB (\lift{\alpha} \prodB (\lift{\beta} \prodB \lift{\gamma}))}{
    \existsB \var{y}, \var{z} \in (\lift{\beta} \prodB \lift{\gamma}) \cdot \\
                        & \hspace{1cm}
    (\var{w} \eqB \var{y} \mapletB \var{z}) \andB
    (\var{x} \mapletB \var{y} \inB Q) \andB
    (\var{x} \mapletB \var{z} \inB R)
    } \quad (\tau = \setB(\alpha \prodB (\beta \prodB \gamma)))
  \end{fitalign}
\end{definition}


\paragraph{Functions and sequences.}
Finally, we introduce additional function-related constructs, since B provides a
rich hierarchy of function spaces.

\begin{definition}[Function spaces][B.Term.tfun, \etc][
    POGReader/Builtins.html\#B.Term.tfun]
  \label{def:b-syntax-functions}
                  \begin{align*}
    A \tfunB B     & \;\triangleq\; \bcomp{\var{f} \inB A \pfunB B}{\forallB \var{x} \inB A \cdot \exists \var{y} \inB B, \var{x} \mapletB \var{y} \inB \var{f}} \tag{total}                               \\
    A \injpfunB B  & \;\triangleq\; \bcomp{\var{f} \inB A \pfunB B}{\forallB \var{x}, \var{y}, \var{z} \inB A \prodB A \prodB B \cdot                                                              \\
                   & \hspace{1cm} (\var{x} \mapletB \var{z} \inB \var{f}) \andB (\var{y} \mapletB \var{z} \inB \var{f}) \impB \var{x} \eqB \var{y}} \tag{partial injective}                                          \\
    A \surjpfunB B & \;\triangleq\; \bcomp{\var{f} \inB A \pfunB B}{\forallB \var{y} \inB B \cdot \existsB \var{x} \inB A \cdot \var{x} \mapletB \var{y} \inB \var{f}} \tag{partial surjective}            \\
    A \injtfunB B  & \;\triangleq\; \bcomp{\var{f} \inB A \tfunB B}{\forallB \var{x}, \var{y} \inB A \prodB A \cdot (\appB \var{f}\ \var{x}) \eqB (\appB \var{f}\ \var{y}) \impB \var{x} \eqB \var{y}} \tag{total injective} \\
    A \surjtfunB B & \;\triangleq\; \bcomp{\var{f} \inB A \tfunB B}{\forallB \var{y} \inB B \cdot \existsB \var{x} \inB A \cdot \var{x} \mapletB \var{y} \inB \var{f}} \tag{total surjective}
  \end{align*}
\end{definition}

Sequences are a special case of functions whose domain is a finite initial
interval of natural numbers.
\begin{definition}[Sequences][B.Term.seq, \etc][
    POGReader/Builtins.html\#B.Term.seq]
  \label{def:b-syntax-sequences}
  Let \(A\) be a B term, and let \(\tau\) be a B type.
  \begin{align*}
    \mathsf{seq}(A)    & \;\triangleq\; \bcomp{\var{s} \inB \mathbb{N} \pfunB A}{\existsB \var{n} \inB \ZintB \cdot \var{s} \inB (\intrangeB{\mathsf{int}\ 1}{\var{n}} \tfunB A)}               \\
    \mathsf{seq}_1(A)  & \;\triangleq\; \bcomp{\var{s} \inB \mathbb{N} \pfunB A}{\left(\existsB \var{n} \inB \ZintB \cdot \var{s} \inB (\intrangeB{\mathsf{int}\ 1}{\var{n}} \tfunB A)\right) \andB \\
                       & \hspace{1cm} \existsB \var{x}, \var{y} \inB \mathbb{N} \prodB A \cdot \appB \var{s}\ \var{x} \eqB \var{y}} \tag{non-empty}                                                        \\
    \mathsf{iseq}(A)   & \;\triangleq\; \mathsf{seq}(A) \capB (\mathbb{N} \injpfunB A) \tag{injective}                                                                               \\
    \mathsf{iseq}_1(A) & \;\triangleq\; \mathsf{seq}_1(A) \capB (\mathbb{N} \injpfunB A) \tag{non-empty injective}                                                                   \\
    \mathsf{perm}(A)   & \;\triangleq\; \mathsf{iseq}(A) \capB (A \surjpfunB A) \tag{permutations}                                                                                   \\
    \mathsf{size}(A)   & \;\triangleq\; \cardB{\domB_{\setB (\intB \prodB \tau)}(A)}
  \end{align*}
\end{definition}

\subsection{Abstract syntax}
\label{subsec:b-abstract-syntax}

The abstract syntax is given by the inductive family
\href{\beerref{B/PHOAS/Basic.lean\#L7}}{\leaninl{B.PHOAS.Term}}, parameterized
by a type \(\mathcal V\) of abstract variables.
The non-binding constructors are structurally the same as in \(\texttt{B.Term}\).
As mentioned in \Cref{sec:phoas}, binding constructors in this setting use
meta-level Lean functions to represent the scope of bound variables:
\[
  \mathsf{binder} : (\mathcal{V} \to \mathsf{Term}\,\mathcal V) \to \mathsf{Term}\,\mathcal V.
\]
We adapt this representation to two specific features of B binders:
\begin{enumerate}[label=(\roman*), nosep]
  \item binders may be polyadic, or \(n\)-ary, since B binders can bind several
        variables simultaneously;
  \item binders require an additional argument for the domain, since they are
        bounded in B.
\end{enumerate}
To avoid nested recursion over lists of variables, we represent the scope of
bound variables functionally, using the type
\[
  \forall n : \mathbb{N},\, \mathsf{Fin}\,n \to \mathcal V,
\]
where \(\mathsf{Fin}\,n\) is the finite ordinal of cardinality \(n\).
This representation is isomorphic to a list-based representation and yields the
following type for binders:
\[
  \mathsf{binder} : \forall n : \mathbb{N},\, \mathsf{Term}\,\mathcal{V} \to ((\mathsf{Fin}\,n \to \mathcal V) \to \mathsf{Term}\,\mathcal V) \to \mathsf{Term}\,\mathcal V.
\]

Overall, the abstract representation of B terms is given by the following
inductive family, parameterized by a type \(\mathcal{V}\) of variables.
To distinguish the two syntactic categories, we reuse the constructor names of
\Cref{def:b-basic-syntax}, but annotate them with a prime symbol.

\begin{definition}[Abstract syntax][B.PHOAS.Term][B/PHOAS/Basic.html\#B.PHOAS.Term]
  \label{def:b-abstract-syntax}
      The constructors for abstract terms are listed below, where \(v\) is a variable in
  \(\mathcal{V}\), \(n\) is an integer, and \(b\) is a Boolean value.

  \begin{bnf}[colspec = {l@{}l@{\kern2pt}c@{\kern2pt}ll}]
    \leaninl{PHOAS.Term} $\,∋\, t, s$ ::=
    | $\var{v}$ && $\mathsf{int}\ n$ && $\mathsf{bool}\ b$ && $\ZintB$ && $\ZboolB$ // Atoms
    | $t \mapletB' s$ // Maplet
    | $t \addB' s$ && $t \subB' s$ && $t \mulB' s$ && $t \leB' s$ // Arithmetic
    | $t \andB' s$ && $\notB' t$ && $t \eqB' s$ && $t \inB' s$ && $\forallB' t \cdot e$ // Logic
    | $\powB'(t)$ && $t \prodB' s$ && $t \cupB' s$ && $t \capB' s$ && $\bcompa{t}{e}$ // Set theory
    | $t \pfunB' s$ && $\appB' t\ s$ && $\lambdaB' t \cdot\,e$ // Functions
  \end{bnf}

  where \(e : (\mathsf{Fin}\,n \to \mathcal{V}) \to \mathsf{Term}\,\mathcal{V}\)
  is the scope of a binder of arity \(n\).
\end{definition}

\begin{remark}
  The arity \(n\) of each binder is implicit, so binders need not be annotated
  with their arity when writing terms.
  It is inferred from the type of the binder scope.
\end{remark}

Substitution can then be defined directly on abstract terms, since the PHOAS
representation delegates \(\alpha\)-conversion and capture-avoidance to Lean's
meta-level binding infrastructure.

\begin{definition}[Substitution][B.PHOAS.Term.subst][B/PHOAS/Basic.html\#B.PHOAS.Term.subst]
  Substitution is defined structurally on abstract terms; the only non-homomorphic
  cases are variables and binders.
  We write \(t[\var{v} \leftarrow e]\) for the capture-avoiding substitution of the
  term \(e\) for the variable \(\var{v}\) in the term \(t\).

  For an arbitrary variable \(\var{w} : \mathcal{V}\), substitution in
  \(\var{w}\) is defined as follows:
  \begin{align*}
    (\var{w}) [\var{v} \leftarrow e]    & \triangleq
    \begin{cases}
      e               & \text{if } \var{w} = \var{v} \\
      \var{w} & \text{otherwise}
    \end{cases}
    \intertext{For a binder \(\int\) with domain \(D\) and scope \(t\), substitution in \(\int_D t\) is defined as follows:}
    \left(\int_D t\right)[\var{v} \leftarrow e] & \triangleq \int_{D[\var{v} \leftarrow e]} \left(\pink{\lambda} x \mapsto (t\ x)[\var{v} \leftarrow e]\right)
    \intertext{where \(\pink{\lambda}\) is the meta-level binder. All remaining clauses are homomorphic. For an \(n\)-ary constructor \(c\) with subterms \(t_1, \ldots, t_n\), substitution is defined as follows:}
    c(t_1, \ldots, t_n)[\var{v} \leftarrow e]   & \triangleq c(t_1[\var{v} \leftarrow e], \ldots, t_n[\var{v} \leftarrow e])
  \end{align*}
\end{definition}

\subsection{Changing representations}
\label{subsec:b-syntax-change-repr}

We now describe the relationship between the concrete and abstract syntaxes, and
how terms can be translated between them.
The relationship is illustrated in \Cref{fig:syntax-representation}.

\begin{figure}[t]
  \centering
  \begin{tikzpicture}[
      node distance=2cm,
      every node/.style={minimum width=4cm, align=center},
    ]
    \node[label=above:{Abstract syntax}, blue] (abstract) at (0, 0) {\(\forallB' \mathbb{N} \cdot \lambda x \mapsto 0 \leB' x\)};
    \node[below of=abstract, label=above:{Concrete syntax}, blue] (concrete) {\(\forallB "x" \inB \mathbb{N} \cdot 0 \leB "x"\)};
    \draw[->, funarr] (concrete.east) to[bend right=30] node[minimum width=2cm, right, label={Abstraction}, blue] {\(\mathsf{String} \to \mathcal{V}\)} (abstract.east);
    \draw[->, funarr] (abstract.west) to[bend right=30] node[minimum width=2cm, left, label={Reification}, blue] {\(\mathcal{V} \hookrightarrow \mathsf{String}\)} (concrete.west);
  \end{tikzpicture}
  \caption{Schematic relationship between abstract and concrete syntax.}
  \label{fig:syntax-representation}
\end{figure}


\paragraph{Reifying abstract terms.}
This direction is conceptually straightforward: one must choose a concrete type
of variable names.
The only subtlety is to recover explicit names from PHOAS binders, which can be
achieved by a suitable fresh-name generation strategy.
This translation is not implemented, because it is unnecessary in our
formalization; we nevertheless sketch the construction for completeness.

The idea is to define a function that, given an abstract term \(t : \mathsf{Term}\,\mathcal{V}\)
over an arbitrary type of variables \(\mathcal{V}\),
produces a reified (concrete) term \( \ulcorner t \urcorner : \mathsf{Term} \).
It is sufficient for \(\mathcal{V}\) to be a countable infinite type---such as
\(\mathsf{String}\) or \(\mathbb{N}\)---to ensure that there are enough
variable names to reify any abstract term with an arbitrary finite number of
bound variables.
Given a function \(\mathsf{fresh}\) that generates fresh variable names,
reification can be defined structurally on abstract terms.
The only non-homomorphic case is that of binders, where fresh variable names
must be generated for the bound variables and substituted in the binder scope.

For this approach to be sound, abstract variables in \(\mathcal{V}\) must be
\emph{injectively embedded} into concrete variable names, written
\(\mathcal{V} \hookrightarrow \mathsf{String}\), so that no variable capture or
identification occurs.
The following example illustrates this reification process.

\begin{example}
  Consider the abstract term \(\lambdaB' \mathbb{N} \prodB \mathbb{N} \cdot \lambda\ x\ y \mapsto x \addB' y\).
  In Lean, one would write the following:
  \begin{leanlisting}
λᴮ' ℕ ×ᴮ ℕ · fun (x : Fin 2 → 𝒱) ↦ .var (x 0) +ᴮ' .var (x 1)
\end{leanlisting}
  Reifying this term with a name generation strategy that produces the names \((\var{a}_i)_{i \in \mathbb{N}}\) yields the following concrete term:
  \[
    \lambdaB \var{a}_0, \var{a}_1 \inB \mathbb{N} \prodB \mathbb{N} \cdot \var{a}_0 \addB \var{a}_1
  \]
  A name generation strategy that produces the same name \(\var{x}\) for both
  variables would instead collapse the two bound variables and yield the
  following non-faithful reification:
  \[
    \lambdaB \var{x}, \var{x} \inB \mathbb{N} \prodB \mathbb{N} \cdot \var{x} \addB \var{x}
  \]
\end{example}


\paragraph{Abstracting concrete terms.}
Conversely, abstraction is used extensively in the subsequent development and
is more delicate.
For a suitably interpreted concrete term, a map \(\Delta\) from concrete
variable names to PHOAS variables provides the interpretation of free variables,
while bound variables are translated into the meta-level function arguments used
by PHOAS.

The abstraction of a concrete term \(t : \mathsf{Term}\) over a map
\(\Delta\), also called a \emph{renaming context}, is defined structurally on
the syntax of \(t\); again, the only non-homomorphic cases are variables and
binders.

\begin{definition}[Abstraction of concrete terms][B.Term.abstract][B/PHOAS/BPHOAS.html\#B.Term.abstract]
  \label{def:b-abstract}
  Let \(t : \mathsf{Term}\) be a concrete term.
  Let \(\mathcal{V}\) be a type of abstract variables, and let
  \(\Delta : \mathsf{String} \rightharpoonup \mathcal{V}\) be a partial map from
  concrete variable names to abstract variables in \(\mathcal{V}\).
  We require only that \(\Delta\) be defined on the free variables of \(t\), so
  that the abstraction of \(t\) over \(\Delta\) is well-defined.

  The abstraction of \(t\) over \(\Delta\), denoted by \(\aterm{t}{\Delta}\),
  is defined structurally on the syntax of \(t\), with variables and binders as
  the only non-homomorphic cases.
  For a variable \(\var{v}\), abstraction is given by lookup in \(\Delta\):
  \begin{align*}
    \aterm{\var{v}}{\Delta}                       & \;\triangleq\; \Delta(\var{v})
    \intertext{For any non-binder \(n\)-ary concrete constructor \(c\) with subterms \(t_1, \dots, t_n\) and abstract counterpart \(c'\), abstraction is defined homomorphically on subterms:}
    \aterm{c(t_1, \dots, t_n)}{\Delta}      & \;\triangleq\; c'(\aterm{t_1}{\Delta}, \dots, \aterm{t_n}{\Delta})
    \intertext{For a concrete binder \(\int\) with abstract counterpart \(\int'\), bound concrete variables \(\vec{\var{v}} = (\var{v_1}, \dots, \var{v_n})\), domain \(D\), and body \(t\), abstraction extends the renaming context with the bound variables and abstracts the body under this extended context:}
    \aterm{\int_{\vec{\var{v}} \inB D} t}{\Delta} & \;\triangleq\; \int_{\aterm{D}{\Delta}}' \left(\pink{\lambda} x_1\ \dots\ x_n \mapsto \aterm{t}{\Delta[\var{v_1} \coloneq x_1, \dots, \var{v_n} \coloneq x_n]}\right)
  \end{align*}
  where \(\Delta[\var{v_1} \coloneq x_1, \dots, \var{v_n} \coloneq x_n]\) denotes the
  extension of \(\Delta\) with the mappings \(\var{v_i} \mapsto x_i\), for
  \(i \in \{1, \dots, n\}\).
\end{definition}
